\section{Introduction}\label{sec:intro}

The input to \sefe{} (\emph{simultaneous embedding with fixed edges} of two
graphs) consists of finite simple planar graphs $G_1=(V,E_1)$ and
$G_2=(V,E_2)$ on the same vertex set. The question is whether there are
planar drawings $\Gamma_1$ of $G_1$ and $\Gamma_2$ of $G_2$ such that every
vertex is drawn at the same point in both and every \emph{common} edge
$e\in E_1\cap E_2$ is drawn as the same curve in both. Edges of
$E_1\setminus E_2$ may cross edges of $E_2\setminus E_1$. We call
$H=(V,E_1\cap E_2)$ the \emph{common graph}, its edges \emph{common}, and all
other edges \emph{private}; $G_1$ is the red graph and $G_2$ the blue graph,
and $U=G_1\cup G_2$ is their union.

Simultaneous embeddings were introduced to visualize several related graphs
on one vertex set; see the survey of Bl\"asius, Kobourov and
Rutter~\cite{BKR13}. Gassner, J\"unger, Percan, Schaefer and
Schulz~\cite{GJPSS06} showed that the problem is NP-complete for three input
graphs, and hardness persists when the union has pathwidth at most
five~\cite{DF26}. For two graphs the complexity has remained open. J\"unger
and Schulz~\cite{JS09} showed that a simultaneous embedding exists if and
only if the two inputs have embeddings that induce the same embedding of
$H$: the same cyclic order of common edges around every vertex and the same
relative positions of the components of $H$.

\paragraph{Known tractable cases.}
\sefe{} is polynomial when $H$ is biconnected~\cite{ADFPR12,HJL13}. For
connected $H$, Fulek and T\'oth~\cite{FT22} gave a polynomial algorithm via
atomic embeddability, and Bl\"asius, Fink and Rutter~\cite{BFR21} a
quadratic one via \emph{Synchronized Planarity}. Schaefer~\cite{Sch13}
showed tractability when every component of $H$ is biconnected or has
maximum degree three. Bl\"asius and Rutter~\cite{BR15} solved the case where
the embedding of every component of $H$ is fixed but the relative positions
are free, which includes disjoint cycles. Bl\"asius, Karrer and
Rutter~\cite{BKR18} gave a cubic algorithm for connected inputs in which
the P-nodes and the simultaneous cutvertices (cutvertices of both inputs but
not of their union) have common degree at most three; they also solve
Schaefer's case in cubic time~\cite[Corollary~1(i)]{BKR18}, which covers
cactus common graphs of maximum degree three.
Parameterized algorithms appear in~\cite{FPR25,DF25,DF26}, and Carmesin and
Turner~\cite{CT25} solve a 2-connected version for temporal sequences of
graphs.

These results leave open the combination of a disconnected common graph,
which forces the drawing to decide which face of one component contains
another, with common vertices of high degree whose cyclic order of common
edges is determined by neither input. Bl\"asius, Karrer and Rutter identify
edge orderings at cutvertices as the central difficulty in this
regime~\cite{BKR18}, and Fink~\cite[Section~6.7]{Fin24} asks whether
Synchronized Planarity can be extended to enforce consistent relative
positions when the common graph is disconnected.
Pfretzschner~\cite[Chapter~8]{Pfr22} carried out such an extension when
both inputs are biconnected and have maximum degree four. He augments
Synchronized Planarity with triples consisting of two edges of a
cycle-basis cycle and an attached edge, whose orientations are linked by
parity constraints, shows how the triples transform under the operations,
and obtains a quadratic algorithm~\cite[Theorem~8.16]{Pfr22}; this covers
cactus common graphs, including disconnected ones, when both inputs are
biconnected of maximum degree four. Our local gap observations and their
transport through PropagatePQ follow the same idea. Unbounded degree
requires three further ingredients: new pipes that carry observations, an
inverse for joining two stars that works component by component, and a
permutation of parallel bonds that preserves all gaps.

\paragraph{Our result.}
A \emph{cactus} is a graph in which every block is a single edge or a simple
cycle; equivalently, two cycles share at most one vertex. We allow a
disconnected common cactus with isolated vertices, bridges, and any number
of cutvertices of any degree. $N$ denotes the number of vertices and edges
of the input.

\begin{theorem}\label{thm:main}
\sefe{} can be decided in $O(N^3)$ time for instances whose common graph is
a cactus. For a positive instance, a pair of compatible combinatorial
embeddings, including the relative positions of all components, is
constructed within the same time bound.
\end{theorem}

Time is counted in operations of a word RAM with words of $\Theta(\log N)$
bits. Combined with a reduction of Bl\"asius, Karrer and
Rutter~\cite{BKR18}, the same bound holds when every component of the common
graph is a cactus or $2$-connected (Corollary~\ref{cor:v4i-mixed}).

\subsection{Overview}\label{sec:intro-overview}
The proof has four steps; each section begins by stating its input and
output.

\paragraph{Normalization.}
Using known reductions we make both inputs connected, replace every common
bridge by a common triangle, and split at cutvertices of the union
(Theorem~\ref{thm:v4-normalize}). In the resulting factors $U$ is
biconnected, both colours are connected, and $H$ is a bridgeless cactus plus
isolated vertices.

\paragraph{Repairing block placements.}
For a common cycle $C$ and a component $Z$ of $U-V(C)$, every simultaneous
embedding draws all of $Z$ on one side of $C$, because $Z$ is connected in
$U$ and the curve of $C$ is common. The repair theorem
(Theorem~\ref{thm:v4-repair}) shows that a weak, block-level form of this
condition suffices. It asks only for embeddings of the blocks of both
colours whose common edge orders at every vertex fit together, and in which
the vertices of the block containing $C$ lie on sides of $C$ that depend
only on their component $Z$. The block embeddings are kept, and only the
attachments at cutvertices are rebuilt. The key fact is that at a vertex on
several common cycles the vectors of side bits form a tree with one edge per
cycle, for every choice of the side bits; so the face structure of the
common graph at that vertex is forced.

\paragraph{Compilation.}
Inside a block, the side of a vertex with respect to a common cycle is the
orientation of a triple of directions at one node of an embedding tree: a
rigid reflection, or the gap of a parallel skeleton between the two
directions carrying the cycle (Lemma~\ref{lem:v4-frontier}(c)). We represent all
blocks by their complete families of embedding trees, which have linear
total size (Lemma~\ref{lem:v4-trees}), glue the red and blue trees at every
vertex along common edges, and obtain an instance of Synchronized Planarity
augmented by marked cycles and parity constraints on local gap observations.
It has linear size (Theorem~\ref{thm:v4-compile}).

\paragraph{Local-gap Synchronized Planarity.}
Rigid reflections are native to Synchronized Planarity, whose Q-vertices
are grouped into consistently oriented cells; gap observations are not. We
run the operations of Bl\"asius, Fink and Rutter, unchanged up to
subdivided edges, and show how each
transforms the observations by explicit parity updates
(Theorem~\ref{thm:eng4-main}). The pipe-free instance at the end is solved
in linear time by reading all observations off the embedding trees of its
blocks (Theorem~\ref{thm:eng4-terminal}).

\paragraph{Where the cactus property is used.}
That common rotations and side bits determine the embedding of $H$ does not
need the cactus property: in every connected plane graph the faces are
distinguished by their sides with respect to the cycles, since two faces are
separated by a minimal edge cut of the dual graph, which is a cycle
(compare~\cite[Theorem~8]{BKR18}). The property that
two common cycles share at most one vertex is used in two places. First,
the common cycles through a vertex form a flower, which gives the forced
tree of side vectors in the repair theorem. Second, a parallel skeleton
displays at most one common cycle (Lemma~\ref{lem:v4-display}), so all gap
observations at a P-vertex refer to a single pair of anchors, and for one
anchor pair every assignment of gap bits is realizable. For a common block
that is neither an edge nor a cycle both properties can fail. Arbitrary
families of gap observations at one vertex need not be describable by
parity constraints (Remark~\ref{rem:eng4-anchors}); this limits a direct
extension of our bookkeeping, but does not show that every extension beyond
cacti must fail.

\subsection{Organization}
Section~\ref{sec:v4-prelim} collects preliminaries and the normalization,
Section~\ref{sec:v4-repair} proves the repair theorem,
Section~\ref{sec:v4-lang} defines the target language,
Section~\ref{sec:v4-compile} recalls SPQR and embedding trees and compiles
normalized instances into it, and
Section~\ref{sec:eng4} solves the compiled instances.
Section~\ref{sec:v4i-conclusion} assembles the running time and discusses
open questions. The appendix proves a variant of a normalization lemma
of Bl\"asius and Rutter~\cite{BR15}.
